\section*{Chapitre \ref{ch:mx:metaorders-latent-liquidity-and-cross-impact} --- Métaordres, liquidité latente et impact croisé}

\begin{solution}{exo:mx:metaorders-latent-liquidity-and-cross-impact:1}
$\sqrt{2\times50/2}=\sqrt{50}=7.07$ ticks.
\end{solution}

\begin{solution}{exo:mx:metaorders-latent-liquidity-and-cross-impact:2}
Les ordres latents reviennent vers le prix plus vite que le métaordre ne les retire, de sorte que celui-ci rencontre toujours un carnet presque plein~: le prix bouge
proportionnellement à ce qu'il a absorbé, comme dans un carnet de profondeur constante.
\end{solution}

\begin{solution}{exo:mx:metaorders-latent-liquidity-and-cross-impact:3}
Avec $I(t)=I_T\sqrt{t/T}$ et un rythme constant, la moyenne vaut $\frac1T\int_0^TI_T\sqrt{t/T}\,dt=\tfrac23I_T$.
\end{solution}

\begin{solution}{exo:mx:metaorders-latent-liquidity-and-cross-impact:4}
La matrice n'est pas symétrique (0.5 contre $-0.1$) et possède un terme diagonal négatif, de sorte qu'un trader pourrait alterner entre les deux actifs en réalisant un profit. La projection
la symétrise et ramène ses valeurs propres négatives à zéro, ce qui donne la matrice symétrique semi-définie positive la plus proche.
\end{solution}

\begin{solution}{exo:mx:metaorders-latent-liquidity-and-cross-impact:5}
En vendant les deux actions, l'\omterm{def:mx:metaorders-latent-liquidity-and-cross-impact:cross}{impact croisé} de chaque vente fait baisser davantage le prix de l'autre, un coût que le modèle d'impact propre ne prend pas en compte (36.0 contre 52.1). Dans le trade de paire,
la vente de la première fait baisser le prix du second pendant qu'on l'achète, un gain que le modèle ne prend pas en compte (36.0 contre 20.0).
\end{solution}

\begin{solution}{exo:mx:metaorders-latent-liquidity-and-cross-impact:6}
La matrice de coût est alors un produit de Kronecker $\Lambda\otimes E$, dont l'inverse est $\Lambda^{-1}\otimes E^{-1}$~: l'optimum vaut $(\Lambda^{-1}\mu)\otimes
E^{-1}\mathbf 1$, soit la même forme temporelle $E^{-1}\mathbf 1$ pour chaque actif, mise à l'échelle de la quantité totale de chacun. Seul le coût dépend des termes croisés.
\end{solution}

\begin{solution}{exo:mx:metaorders-latent-liquidity-and-cross-impact:7}
Pour une exécution de 100 secondes, l'exposant reste à 1.0 jusqu'à $Q=100$ et ne tombe à 0.94 et 0.75 que pour 300 et 1\,000~: le point de bascule se déplace d'un facteur dix, vers
environ $LDT=100$, car le carnet dispose de dix fois plus de temps pour se reconstituer.
\end{solution}

\begin{solution}{exo:mx:metaorders-latent-liquidity-and-cross-impact:8}
Le prix qui repasse sous le prix moyen d'exécution après la fin est ce que prédit le prix équitable avec un impact transitoire, et ne prouve la perte de personne~; le prix après
la fin continue de bouger (dans le carnet latent, il continue de baisser). Qui a perdu dépend de la durée de détention des exécutants et du moment où ils ont clôturé~; il faut mesurer leurs mark-out
à leurs propres horizons.
\end{solution}

\begin{solution}{pb:mx:metaorders-latent-liquidity-and-cross-impact:1}
\textbf{1.} Fractionner l'ordre avec un impact décroissant, une information et un prix équitable, et un carnet latent qui disparaît au prix.
\textbf{2.} Intentions de trading non placées dans le carnet~; une densité latente $L|x-p|$.
\textbf{3.} $\int_0^ILx\,dx=LI^2/2=Q$.
\textbf{4.} 0.05, 0.18, 0.52, 1.74, 4.9, 12.2, 23.3 et 44.0 ticks~; exposant 1.0 jusqu'à $Q=10$, puis 0.94, 0.76, 0.59, 0.53~: le point de bascule proche de $LDT=10$.
\textbf{5.} Le prix moyen payé est égal au prix après exécution complète.
\textbf{6.} Les deux tiers du pic.
\textbf{7.} 1.9 seconde après la fin pour $Q=10$ (2.3 pour 300)~; il continue de baisser, jusqu'à 0.16 (0.27) du pic après 100 secondes.
\textbf{8.} Les ordres latents ne perdent jamais la mémoire dans cette version~: le carnet se reconstitue autour de l'ancien prix~; le dépôt et l'annulation d'ordres latents
ancreraient une partie du mouvement.
\textbf{9.} $r_t=\Lambda q_t+\text{bruit}$.
\textbf{10.} Un déséquilibre sur une action ou un secteur déplace les rendements des autres, de façon persistante, souvent négativement~; les transactions médiatisent une grande partie de la covariance des rendements,
la liquidité partageant la structure sectorielle des corrélations.
\textbf{11.} $\begin{pmatrix}0.96&0.42\\0.42&0.81\end{pmatrix}$~; 4.5~\% et 5.3~\% de la variance des rendements au-delà du flux propre de chaque action.
\textbf{12.} \omterm{def:mx:metaorders-latent-liquidity-and-cross-impact:cross}{Impact croisé} linéaire, impair et symétrique, pour des noyaux bornés, sinon des cycles à travers les actifs seraient rentables.
\textbf{13.} 52.1 contre 36.0.
\textbf{14.} 20.0 contre 36.0.
\textbf{15.} 64.6 pour la vente et 38.4 pour la paire~: l'échéancier conjoint économise 19~\% et 48~\%.
\textbf{16.} Lorsque l'\omterm{def:mx:metaorders-latent-liquidity-and-cross-impact:cross}{impact croisé} décroît à un rythme différent de l'impact propre~; avec une décroissance croisée plus lente, le trade de paire économise 3.7~\% de plus.
\textbf{17.} \emph{Résultat notable}~: le flux de l'autre action explique 4.5~\% et 5.3~\% de la variance des rendements de chacune au-delà de son propre flux~; négocier les deux jambes ensemble
plutôt que l'une après l'autre économise 19~\% du coût de la vente et 48~\% du trade de paire, et un modèle d'impact propre fausse le coût de $-31~\%$ et de $+80~\%$.
\textbf{18.} Les deux actions ensemble sur tout l'horizon, avec une estimation du coût qui inclut l'\omterm{def:mx:metaorders-latent-liquidity-and-cross-impact:cross}{impact croisé}.
\textbf{19.} Les deux jambes simultanément, de sorte que l'\omterm{def:mx:metaorders-latent-liquidity-and-cross-impact:cross}{impact croisé} de chaque jambe aide l'autre.
\textbf{20.} Son cours bouge avec nos transactions sur l'action que nous vendons, et il en va de même pour nos coûts et nos valorisations sur celle-ci.
\end{solution}

\begin{solution}{iq:mx:metaorders-latent-liquidity-and-cross-impact:1}
La majeure partie de la liquidité est latente, et sa densité croît linéairement en s'éloignant du prix. Déplacer le prix de $I$ consomme $LI^2/2$ de celle-ci. Ainsi, un métaordre de $Q$
déplace le prix de $\sqrt{2Q/L}$ lorsque le carnet ne peut pas se reconstituer à temps.
\iqlookfor{Carnet latent~; l'intégrale~; le régime rapide.}
\end{solution}

\begin{solution}{iq:mx:metaorders-latent-liquidity-and-cross-impact:2}
Le prix moyen payé est égal au prix après exécution complète~; avec un impact en racine carrée, le prix devrait se stabiliser vers les deux tiers du pic.
\iqlookfor{La condition~; la conséquence des deux tiers.}
\end{solution}

\begin{solution}{iq:mx:metaorders-latent-liquidity-and-cross-impact:3}
Des positions sur des noms corrélés, par l'\omterm{def:mx:metaorders-latent-liquidity-and-cross-impact:cross}{impact croisé} sur leurs cours~: les valorisations et les ventes ultérieures en pâtissent. Estimer l'\omterm{def:mx:metaorders-latent-liquidity-and-cross-impact:cross}{impact croisé} sur mes propres transactions ou sur le flux d'ordres du marché,
en régressant les rendements de chaque nom sur les flux des autres.
\iqlookfor{\omterm{def:mx:metaorders-latent-liquidity-and-cross-impact:cross}{Impact croisé}~; estimation sur les flux.}
\end{solution}

\begin{solution}{iq:mx:metaorders-latent-liquidity-and-cross-impact:4}
Régresser les rendements par intervalle sur le vecteur des déséquilibres du flux d'ordres, avec assez d'intervalles et des noms regroupés par secteur~; imposer la symétrie et la semi-définie positivité
(projeter l'estimation), et envisager une décomposition en valeurs propres de la covariance des flux pour réduire le bruit.
\iqlookfor{Régression sur les flux~; symétrie et semi-définie positivité~; réduction de dimension.}
\end{solution}

\begin{solution}{iq:mx:metaorders-latent-liquidity-and-cross-impact:5}
Un schéma explicite exige $D\,dt/dx^2\le\tfrac12$ pour la stabilité~; la frontière (le prix) est trouvée à chaque pas comme le zéro de la densité, interpolé
entre les cellules~; un schéma implicite supprime la limite sur le pas, au prix d'une résolution linéaire.
\iqlookfor{La condition de stabilité~; le suivi de la frontière.}
\end{solution}

\begin{solution}{iq:mx:metaorders-latent-liquidity-and-cross-impact:6}
Avec un \omterm{def:mx:metaorders-latent-liquidity-and-cross-impact:cross}{impact croisé} positif, vendre une jambe fait baisser le prix de l'autre tandis que l'acheter le fait monter~: les termes croisés des jambes se compensent. Négocier les deux jambes simultanément et
de façon équilibrée, et estimer le coût en incluant les termes croisés.
\iqlookfor{\omterm{def:mx:metaorders-latent-liquidity-and-cross-impact:cross}{Impact croisé} compensateur~; exécution simultanée.}
\end{solution}
